\section{Results}
\label{sec:results}

\subsection{MOFA captures a shared multi-omics latent space}

Of the 20 initialised factors, 13 were retained after ARD filtering
($R^2_{\max} \geq 0.01$).
Factor~1 is by far the dominant component, explaining 13.0\% of transcriptomic,
33.6\% of proteomic, and 32.1\% of phosphoproteomic variance
(Figure~\ref{fig:var_explained}).
The remaining 12 factors each explain less than 8.1\% of variance in any single view,
with most falling below 3\%, indicating that variance is heavily concentrated in Factor~1.

Factor~1 perfectly separates tumour from normal tissue: a paired Wilcoxon test on
$\Delta Z$ yielded $p = 1.8 \times 10^{-18}$ with rank-biserial $r = 1.0$, meaning
all 102 patients shifted in the same direction along Factor~1 during tumourigenesis.
No other factor reached statistical significance.
The PCA trajectory plot in Figure~\ref{fig:traj_pca} makes this structure visible:
PC1 is dominated by Factor~1, and the 102 per-patient arrows are predominantly
parallel and oriented along this axis.
The consistency of the Factor~1 trajectory confirms that MOFA has learned a
biologically coherent shared space in which tumourigenesis appears as a
reproducible displacement.

\begin{figure}[t]
    \centering
    \includegraphics[width=0.55\textwidth]{figures/variance_explained_heatmap.pdf}
    \caption{Fraction of variance explained by each retained MOFA+ factor per
    omics view. Factor~1 explains substantially more variance than any other factor
    across all three modalities.}
    \label{fig:var_explained}
\end{figure}

\begin{figure}[t]
    \centering
    \includegraphics[width=0.85\textwidth]{figures/trajectory_arrows_pca.pdf}
    \caption{Per-patient tumourigenesis trajectories in the PCA of MOFA+ factor
    scores (all 204 samples). Each arrow connects the normal (open circle) to the
    tumour (filled circle) sample of one patient. The dominant direction along PC1
    reflects the near-universal Factor~1 shift.}
    \label{fig:traj_pca}
\end{figure}

\subsection{$\Delta Z$ vectors encode patient-level transformation heterogeneity}
\label{sec:delta}

The $\Delta Z$ matrix for 102 patients shows that Factor~1 drives a large, uniform
component of tumourigenesis: the rank-biserial correlation of $r = 1.0$ implies that
the tumour-minus-normal displacement is in the same direction for every patient along
this axis.
In contrast, the remaining 12 factors show substantial variation across patients
(per-factor distributions are shown in Appendix Figure~\ref{fig:delta_distributions}),
suggesting that inter-patient heterogeneity in transformation is concentrated in the
dimensions beyond the dominant tumour-normal axis.

Transformation magnitude (L2 norm of $\Delta Z$) varied across patients but showed
no significant association with tumour stage
(Kruskal-Wallis $H = 3.40$, $p = 0.183$; Appendix Figure~\ref{fig:mag_by_stage}).
The absence of a stage-magnitude relationship indicates that transformation extent
in latent space does not simply reflect disease progression as captured by clinical
staging.

\subsection{Clustering identifies three reproducible transformation subtypes}

We swept $k \in \{2, 3, 4, 5, 6\}$ on L2-normalised $\Delta Z$.
The silhouette score peaked at $k = 2$ (0.184) and remained nearly identical at
$k = 3$ (0.182), after which it declined monotonically (Table~\ref{tab:model_sel}).
GMM BIC increased continuously with $k$, while AIC continued to decrease, indicating
that the penalty for additional components outweighs the likelihood gain and that the
data do not support strongly separated Gaussian clusters (Figure~\ref{fig:model_sel}).
We selected $k = 3$ over $k = 2$ based on the nearly indistinguishable silhouette
scores, the more balanced partition (37, 25, and 40 patients versus 77 and 25 for
$k = 2$), and the fact that $k = 3$ yields a biologically richer description of
inter-patient heterogeneity (see \S\ref{sec:interpret}).

\begin{table}[t]
    \centering
    \caption{Clustering quality metrics across $k$ on L2-normalised $\Delta Z$.
    Silhouette scores are means over all patients; GMM AIC and BIC are computed
    independently for each $k$ as descriptive metrics.}
    \label{tab:model_sel}
    \begin{tabular}{rrrr}
        \toprule
        $k$ & Silhouette & GMM AIC & GMM BIC \\
        \midrule
        2 & 0.184 & $-$1298.2 & $-$749.6 \\
        3 & 0.182 & $-$1303.5 & $-$479.3 \\
        4 & 0.169 & $-$1530.2 & $-$430.3 \\
        5 & 0.165 & $-$1574.7 & $-$199.2 \\
        6 & 0.157 & $-$1842.4 & $-$191.3 \\
        \bottomrule
    \end{tabular}
\end{table}

\begin{figure}[t]
    \begin{subfigure}[t]{0.56\textwidth}
        \centering
        \includegraphics[width=\textwidth]{figures/model_selection.pdf}
        \caption{Silhouette score (left) and GMM AIC/BIC (right) versus $k$.
        The silhouette plateau at $k = 2$--$3$ and the increasing BIC are consistent
        with weak cluster separation on a continuous manifold.}
        \label{fig:model_sel}
    \end{subfigure}
    \hfill
    \begin{subfigure}[t]{0.40\textwidth}
        \centering
        \includegraphics[width=\textwidth]{figures/cross_repr_ari.pdf}
        \caption{Adjusted Rand index between the primary partition
        ($k = 3$, L2-normalised $\Delta Z$) and three alternative representations.
        All ARI values exceed 0.79.}
        \label{fig:ari}
    \end{subfigure}
    \caption{Clustering model selection (a) and cross-representation robustness (b).}
    \label{fig:clustering_diag}
\end{figure}

The three-subtype partition was robust across representational choices: ARI against
L2-normalised $\Delta Z$ without Factor~1 was 0.82, against raw $\Delta Z$ was 0.80,
and against 2D PCA of the primary representation was 0.80
(Figure~\ref{fig:ari}).
Pairwise centroid distances in L2-normalised space ranged from 0.56 (subtypes 0 and 2)
to 0.83 (subtypes 0 and 1), exceeding the mean within-cluster radius of
0.48--0.54 for all pairs.

The UMAP projection in Figure~\ref{fig:cluster_umap} and the $\Delta Z$ heatmap in
Figure~\ref{fig:delta_heatmap} illustrate the structure of the three subtypes.
In the UMAP, the subtypes occupy partially overlapping regions consistent with a
continuous manifold rather than discrete blobs, in agreement with the silhouette
scores.
The heatmap confirms that all subtypes share a strongly negative Factor~1 component
(universal tumourigenesis shift) but differ in their profiles across Factors 2--13,
motivating the biological characterisation below.

\begin{figure}[t]
    \centering
    \includegraphics[width=0.55\textwidth]{figures/cluster_umap.pdf}
    \caption{UMAP of L2-normalised $\Delta Z$ coloured by transformation subtype.
    Clustering was performed in the full 13-dimensional space; the UMAP is for
    visualisation only. Partial overlap between subtypes is consistent with the
    low silhouette scores.}
    \label{fig:cluster_umap}
\end{figure}

\begin{figure}[t]
    \centering
    \includegraphics[width=\textwidth]{figures/delta_heatmap.pdf}
    \caption{Heatmap of raw $\Delta Z$ values (patients $\times$ factors) sorted by
    transformation subtype. All subtypes share a uniformly negative Factor~1
    component; inter-subtype differences are concentrated in Factors 2--13.}
    \label{fig:delta_heatmap}
\end{figure}

\subsection{Transformation subtypes reflect distinct molecular programmes}
\label{sec:interpret}

\paragraph{Factor interpretation.}
Preranked GSEA on RNA factor loadings yielded significant results (FDR $< 0.10$) for
8 of the 13 retained factors; Factors~1, 2, 4, 6, and 9 returned no significant
enrichment at this threshold (Figure~\ref{fig:gsea_factor}).
The most consistently enriched categories were cytokine-cytokine receptor interaction
(Factors 3, 5, 13), xenobiotic metabolism and chemical carcinogenesis (Factors 10
and 11), and downstream KRAS signalling suppression (Factors 7 and 12).
Factor~1, despite its strong tissue-separation signal, showed no significant pathway
enrichment, suggesting that it captures a broad transcriptional state shift that is
not well-represented by the curated gene sets used here.

\begin{figure}[t]
    \centering
    \includegraphics[width=\textwidth]{figures/gsea_factor_heatmap.pdf}
    \caption{Preranked GSEA results for retained MOFA+ factors using RNA loading
    weights. Cells show NES; only pathways significant in at least one factor
    (FDR $< 0.10$) are displayed. Positive NES indicates enrichment at the positive
    loading pole; negative NES at the negative pole.}
    \label{fig:gsea_factor}
\end{figure}

\paragraph{Cluster interpretation.}
One-vs-rest $\Delta$Expression GSEA revealed that all three subtypes are
transcriptionally distinct (Figure~\ref{fig:gsea_cluster}).
Subtype~1 ($n = 25$) showed the strongest activation of proliferative programmes:
E2F targets (NES $= 1.80$, FDR $< 0.001$) and G2-M checkpoint
(NES $= 1.75$, FDR $= 0.002$) were among the most significantly enriched pathways,
indicating that these patients undergo substantially greater tumour-vs-normal
upregulation of cell cycle machinery than patients in the other two subtypes.
Subtype~0 ($n = 37$) displayed the inverse pattern: the same E2F, G2-M, and MYC
target gene sets were negatively enriched (NES $\leq -1.61$, FDR $\leq 0.04$),
reflecting relatively suppressed proliferative transformation compared to subtypes~1
and 2.
Subtype~2 ($n = 40$) showed the strongest interferon response: interferon-$\alpha$
(NES $= 1.79$, FDR $= 0.008$) and interferon-$\gamma$ (NES $= 1.77$, FDR $= 0.007$)
were the top enriched pathways.
E2F targets were also enriched in subtype~2 (NES $= 1.68$, FDR $= 0.035$), in
contrast to their depletion in subtype~0.
The three subtypes therefore differ systematically in their transcriptional
reprogramming: subtype~1 is proliferation-dominant, subtype~0 shows relatively
suppressed cell-cycle activation, and subtype~2 combines immune activation with a
proliferative component.

\begin{figure}[t]
    \centering
    \includegraphics[width=\textwidth]{figures/gsea_cluster_heatmap.pdf}
    \caption{One-vs-rest $\Delta$Expression GSEA results for the three transformation
    subtypes. Cells show NES; only pathways significant in at least one subtype
    (FDR $< 0.10$) are displayed. Negative NES indicates that the pathway is
    relatively suppressed in the focal subtype compared with the other patients.}
    \label{fig:gsea_cluster}
\end{figure}

\subsection{Preliminary clinical associations}

No significant difference in overall survival was observed across the three subtypes
(omnibus log-rank $p = 0.632$; Figure~\ref{fig:km}). Hence, pairwise comparisons were not
performed.
With $n = 102$ patients, the study is insufficiently powered to detect modest
survival differences.

Driver mutation enrichment analysis revealed a significant association between
subtype and \textit{EGFR} mutation status (Figure~\ref{fig:mut}).
\textit{EGFR} mutations were enriched in subtype~0 (odds ratio 3.18,
BH-adjusted $p = 0.043$) and depleted in subtype~1 (odds ratio 0.22,
BH-adjusted $p = 0.043$).
ALK fusions were strongly depleted in subtype~0 (odds ratio 0.12,
BH-adjusted $p = 0.016$), consistent with the mutual exclusivity of
\textit{EGFR} mutations and ALK rearrangements as driver events.
\textit{KRAS} and \textit{TP53} mutations showed no significant cluster association
after Benjamini-Hochberg correction.
The co-occurrence of EGFR enrichment, suppressed proliferative transformation, and
relatively lower cell-cycle upregulation in subtype~0 is consistent with the
biology of EGFR-driven LUAD, which tends towards a more differentiated, TRU-like
transcriptional state.

\begin{figure}[t]
    \centering
    \includegraphics[width=0.7\textwidth]{figures/km_curves.pdf}
    \caption{Kaplan-Meier overall survival curves for the three transformation
    subtypes. Shading shows 95\% confidence intervals; tick marks indicate
    censored observations. Omnibus log-rank $p = 0.632$.}
    \label{fig:km}
\end{figure}

\begin{figure}[t]
    \centering
    \includegraphics[width=0.72\textwidth]{figures/mutation_enrichment_heatmap.pdf}
    \caption{Driver mutation enrichment across transformation subtypes. Cell colour
    encodes $\log_2$ odds ratio; asterisks indicate BH-adjusted $p < 0.05$ ($*$)
    or $p < 0.01$ ($**$).}
    \label{fig:mut}
\end{figure}
